\capitulo{2}{Objetivos del proyecto}

Este apartado explica de forma precisa y concisa cuales son los objetivos que se persiguen con la realización del proyecto. Se puede distinguir entre los objetivos marcados por los requisitos del software a construir y los objetivos de carácter técnico que plantea a la hora de llevar a la práctica el proyecto.

El objetivo que persigo con este trabajo es un reto importante ya que lo que se busca es unir dos mundos, virtual y presencial, mediante la utilización de tecnología existente.

Se trata de diseñar un laboratorio virtual que permita a un usuario programar un robot e incluso ver como se construye el mismo para hacer un montaje presencial del mismo. Además, la apuesta mas arriesgada es la de ser capaz de ensamblar un robot a través del ordenador o dispositivo móvil.

Además supone adquirir formación en el motor de videojuegos Unity que se va a utilizar y para ello hay que diseñar una interface que integre el entorno de programación de scracth, herramienta open source, con un entorno virtual donde el alumno pueda ver como es todo el proceso desde la creación del robot hasta la programación y puesta final.

Unity tiene muchas capacidades que me van a permitir aprender todo lo relacionado con este programa que puede hacer realidad una necesidad no sin los inconvenientes que van a existir.

En el aspecto técnico es necesario dominar los conceptos de diseño que necesita el motor para llevar a cabo el proceso de creación mediante la programación con lenguaje C# de las distintas clases que van a permitir operar con todos los elementos de la aplicación.

Existe una problemática inicial, aunque scratch sea open source, que impide integrar directamente los bloques y su lógica en Unity porque este no maneja JS y CSS.

Los bloques de scratch además se adaptan según van conteniendo parámetros, operadores o sensores con lo que se hace necesario el aprendizaje de conceptos relacionados con la UI de Unity para el correcto manejo de los mismos.

Por otro lado, la programación de robots de la marca Lego se realiza utilizando sus propios entornos aunque también las versiones Wedo y EV3 se pueden programar con scracth vinculando los distintos elementos.

Dichos robots que son costosos para cualquier persona, al ser desarrollados de manera virtual no tendrían un coste tan elevado y podría ser asumido por cualquier usuario usando una aplicación para que se lleven a cabo pruebas con los mismos usando las físicas de Unity el cual ya las tiene configuradas e imitan el mundo real.

En el campo de la programación no me va a suponer adquirir conocimientos ya que los he ido recibiendo durante la titulación pero la parte técnica que más me va recursos me va a consumir es el conocimiento de la parte de diseño de Unity para que la escena funcione adecuadamente.




